%% LyX 2.4.4 created this file.  For more info, see https://www.lyx.org/.
%% Do not edit unless you really know what you are doing.
\documentclass[english]{article}
\usepackage[T1]{fontenc}
\usepackage[latin9]{inputenc}
\usepackage{enumitem}
\usepackage{amsmath}
\usepackage{amsthm}
\usepackage{amssymb}
\usepackage{geometry}
\geometry{verbose,tmargin=1in,bmargin=1in,lmargin=1in,rmargin=1in}

\makeatletter
%%%%%%%%%%%%%%%%%%%%%%%%%%%%%% Textclass specific LaTeX commands.
\numberwithin{equation}{section}
\numberwithin{figure}{section}
\newlength{\lyxlabelwidth}      % auxiliary length 

%%%%%%%%%%%%%%%%%%%%%%%%%%%%%% User specified LaTeX commands.
\usepackage{fancyhdr}
\usepackage{lastpage}
\pagestyle{fancy}
\usepackage{enumitem}
\usepackage{ifthen}
%sets math fonts to be more readable
\everymath=\expandafter{\the\everymath\displaystyle}
%\DeclareMathSizes{10}{10}{10}{10}
\usepackage{multicol}

\usepackage{advdate}    % Advancing/saving dates
\usepackage[dayofweek]{datetime}   % Dates formatting
\usepackage{datenumber} % Counters for dates
% Added by lyx2lyx
\setlength{\parskip}{\smallskipamount}
\setlength{\parindent}{0pt}

\makeatother

\usepackage{babel}
\begin{document}
\renewcommand{\author}{Dr.\,James Ashe}

\newcommand{\classNum}{335}

\newcommand{\classSec}{A}

\newcommand{\nameType}{}

\newcommand{\examType}{Homework}\usdate

\newcommand{\Date}{Due date: \formatdate{20}{2}{2014}} %%\AdvanceDate[12] \today

%\newdate{my_date}{\day}{\month}{\year}%   
%\AdvanceDate[-5] 
%\displaydate{my_date}

%%First page's header%%%%%%%%%%%%%%%%%%%%%%
\fancypagestyle{firststyle} %%header settings for first pagr
{
	\fancyhead[L]{MTH \classNum{}(\classSec{}) \author{}\\
	\textbf{\examType{}} \Date{}}
	\fancyhead[C]{\textbf{\nameType}}
	\setlength{\headsep}{14pt}
}
%%%%%%%%%%%%%%%%%%%%%%%%%%%%%%%%%%%%%%%%%%

%%Next page's header%%%%%%%%%%%%%%%%%%%%%%%
%\setlist{nolistsep} %eliminates space between list items
\lhead{MTH \classNum{}(\classSec{})} 
\chead{\textbf{\nameType}} 
\cfoot{\thepage{}~of~\pageref{LastPage}} 
\setlength{\headsep}{5pt} 
%%%%%%%%%%%%%%%%%%%%%%%%%%%%%%%%%%%%%%%%%%%

%%Various settings; see the preamble for other settings%%%%%
%%\setenumerate[0]{leftmargin=12pt,itemindent=0pt,labelwidth=0pt}
\renewcommand{\labelenumi}{\textbf{\arabic{enumi}.}}
\renewcommand{\labelenumii}{\textbf{\alph{enumii}.}}
\thispagestyle{firststyle}
%%%%%%%%%%%%%%%%%%%%%%%%%%%%%%%%%%%%%%%%%%

\textbf{Homework Problems. }\emph{Solutions should be written/typed
in numerical order and clearly labeled. Unsolved problems should be
included but left blank.}
\begin{enumerate}
\item Prove or disprove the following.

\begin{enumerate}
\item A constant function is not one-to-one if the domain has more than
one element.
\end{enumerate}
\begin{proof}
[True.]Let $f\colon A\rightarrow B$ be a constant function and $|A|>1$.
Since $|A|>1$ it has at least two elements, say $a_{1}$ and $a_{2}$.
$f$ is constant so then $f(a_{1})=f(a_{2})=b$ for some $b\in B$.
Thus $f$ is not one-to-one.

Note that if the domain has only one element then this statement is
false. Let $A=\{a\}$ and $f:A\rightarrow B$ and $f(a)=b$ for all
$a\in A$. If $f(a_{1})=f(a_{2})$ then $a_{1}=a_{2}$ because $A$
only has one element. So $f$ is constant and also one-to-one. \emph{In
general constant functions are not 1-1, i.e., the statement ``constant
functions are 1-1'' is false.}
\end{proof}
\begin{enumerate}[resume]
\item The identity function is one-to-one.
\end{enumerate}
\begin{proof}
[True.] Let $f:A\rightarrow A$ be the identity function. If $f(a_{1})=f(a_{2})$
then $f(a_{1})=a_{1}=a_{2}=f(a_{2})$ because $f(a)=a$ for any $a\in A$.
Thus $f$ is one-to-one. 
\end{proof}
\begin{enumerate}[resume]
\item The identity function is onto.
\end{enumerate}
\begin{proof}
[True.]Let $f:A\rightarrow A$ be the identity function and $a\in A$.
$f(a)=a$ so then $f$ is onto.
\end{proof}
\end{enumerate}
\vfill
\begin{enumerate}[resume]
\item Let $g:A\rightarrow B$ and $f:B\rightarrow C$. Prove that if $f\circ g:A\rightarrow C$
is one-to-one and $g$ is onto then $f$ is also one-to-one. 

\begin{proof}
Let $b_{1},b_{2}\in B$ such that $b_{1}\neq b_{2}$. Because $g$
is onto there is a $a_{1},a_{2}\in A$ such that $g(a_{1})=b_{1}$
and $g(a_{2})=b_{2}$. So then $f(b_{1})\neq f(b_{2})$ since $f(g(a_{1}))\neq f(g(a_{2}))$. 
\end{proof}
\end{enumerate}
\vfill
\begin{enumerate}[resume]
\item Suppose that $f:A\rightarrow B$ is a one-to-one function. Prove or
disprove that if $f^{-1}:f(B)\rightarrow A$ is a function then $f^{-1}$
is also one-to-one.

\begin{proof}
[True.] Recall: It is shown in the notes that for a function $g$,
the inverse of $g$ is a function if and only if $g$ is $1-1$. Since
$f^{-1}$ is a function and $(f^{-1})^{-1}=f$ is also a function,
it follows that $f^{-1}$ is $1-1$. 
\end{proof}
Alternatively,
\begin{proof}
Suppose to the contrary that $f^{-1}$ is \emph{not }$1-1$. So then
there is $b_{1},b_{2}\in B$ such that $b_{1}\neq b_{2}$ and $f^{-1}(b_{1})=f^{-1}(b_{2})=a$
for some $a\in A$. Thus $f(f^{-1}(b_{1}))=b_{1}\neq b_{2}=f(f^{-1}(b_{2}))$
which contradicts the assumption that $f$ is a fucntion. 
\end{proof}
\end{enumerate}
\vfill
\begin{enumerate}[resume]
\item Let $f:A\rightarrow B$ be a one-to-one function. Show that $\left|A\right|\leq\left|B\right|$.

\begin{proof}
Let $\left|A\right|$ be the number of items and $\left|B\right|$
be the number of boxes. By the pigeon hole theorem, if $\left|A\right|>\left|B\right|$
then $f$ would map (or place) at least $2$ items to the same box
and thus not be $1-1$. Thus $f$ must be $1-1$.
\end{proof}
\end{enumerate}
\vfill\newpage
\begin{enumerate}[resume]
\item Let $f:A\rightarrow B$ and $g:B\rightarrow C$ be functions. Prove
that $g\circ f:A\rightarrow C$ is a function. 

\begin{proof}
$f$ is a function implies that for any $a\in A$ there is a unique
$b\in B$ such that $f(a)=b$. Also, $g$ is a function implies that
for any $b\in B$ there is a unique $c\in C$ such that $g(b)=c$.
Thus for any for any $a\in A$ there is a unique $c\in C$ such that
$g(f(a))=g(b)=c$. So $g\circ f$ is a function.
\end{proof}
\end{enumerate}
\vfill
\begin{enumerate}[resume]
\item (This problem was dropped) Let $f:A\rightarrow A$, $g:A\rightarrow A$,
and $h:A\rightarrow A$; prove that $\left(h\circ g\right)\circ f=h\circ\left(g\circ f\right)$.

\begin{proof}
Let $a_{0}\in A$. Since $f,g,h$ are functions there are $a_{1},a_{2},a_{3}\in A$
such that, $f(a_{0})=a_{1}$, $g(a_{1})=a_{2}$, and $h(a_{2})=a_{3}$.
$\left(\left(h\circ g\right)\circ f\right)(a_{0})=h(g(f(a_{0})))=h(g(a_{1}))=h(a_{2})=a_{3}$.
Also, $\left(h\circ\left(g\circ f\right)\right)(a_{0})=h((g\circ f)(a_{0}))=h(g(f(a_{0})))=h(g(a_{1}))=h(a_{2})=a_{3}$.
Thus $\left(h\circ g\right)\circ f=h\circ\left(g\circ f\right)$.
\end{proof}
\end{enumerate}
\vfill
\begin{enumerate}[resume]
\item Prove that if $f:A\rightarrow B$ is onto and $g:B\rightarrow C$
is onto, then $\left(g\circ f\right):A\rightarrow C$ is onto. 

\begin{proof}
Let $c\in C$. $g$ is onto so then there exists a\textbf{ $b\in B$
}such that $g(b)=c$. Since $f$ is onto there also exists an $a\in A$
such that $f(a)=b$. So then $g(f(a))=g(b)=c$; thus $g\circ f$ is
onto.
\end{proof}
\end{enumerate}
\vfill
\begin{enumerate}[resume]
\item Each of the following statements, $P(x,y)$, define a relation $R$
in the natural numbers, $\mathbb{N}$, i.e., $R=(\mathbb{N},\mbox{\ensuremath{\mathbb{N}}},P(x,y))$.
Show whether each of the following relations is reflexive. 

\begin{enumerate}
\item $P(x,y)=$``$x$ is less than or equal to $y$''
\end{enumerate}
\begin{proof}
[Reflexive:] Let $x\in\mathbb{N}$. $x=x\Rightarrow$$x\leq x$ so
then $R$ is reflexive.
\end{proof}
\begin{enumerate}[resume]
\item $P(x,y)=$``$x$ divides $y$''
\end{enumerate}
\begin{proof}
[Reflexive:] Let $x\in\mathbb{N}$. $x=1\cdot x$ so then $x|x$.
Thus $R$ is reflexive.
\end{proof}
\begin{enumerate}[resume]
\item $P(x,y)=$``$x+y=10$''
\end{enumerate}
\begin{proof}
[Not Reflexive:] Consider the natural number $4$; $4+4\neq10$.
Thus $R$ is not reflexive. 
\end{proof}
\begin{enumerate}[resume]
\item $P(x,y)=$ ``$x$ and $y$ are relatively prime.''
\end{enumerate}
\begin{proof}
[Not Reflexive:] Consider the natural number $8$; $\gcd(8,8)=8\neq1$
since $8|8$. Thus $R$ is not reflexive.
\end{proof}
\end{enumerate}
\vfill
\end{document}
